% Source: https://artofproblemsolving.com/community/c3h1402095_a_problem_from_an_mc_mock_target
% Stable identity: aops-topic-1402095-q01
Let $AB=a$ and $CD=b$ with $a>b$. Triangles $EAB$ and $ECD$ are similar (as $AB\parallel CD$), so $\dfrac{CE}{EA}=\dfrac{b}{a}$ and $\dfrac{CE}{CA}=\dfrac{b}{a+b}$.

\begin{asy}
usepackage("amsmath");
size(6.5cm);
defaultpen(fontsize(10pt));
pair A=(0,0),B=(24,0),C=(18,8),D=(6,8);
pair Ep=extension(A,C,B,D); pair F=extension(Ep,Ep+(1,0),B,C);
draw(A--B--C--D--cycle,linewidth(0.9));
draw(A--C,gray); draw(B--D,gray);
draw(Ep--F,red+linewidth(0.9));
dot(A);dot(B);dot(C);dot(D);dot(Ep);dot(F,red);
label("$A$",A,SW); label("$B$",B,SE); label("$C$",C,NE); label("$D$",D,NW); label("$E$",Ep,S); label("$F$",F,E,red);
label("$a$",(A+B)/2,S); label("$b$",(C+D)/2,N); label("$8$",(Ep+F)/2,N,red);
label("\mbox{}",(min(currentpicture,user=true)+max(currentpicture,user=true))/2,black); // keeps the text after the figure black when LaTeX typesets the labels inline
\end{asy}

Since $EF\parallel AB$, triangle $CEF$ is similar to triangle $CAB$, so
\[
EF=AB\cdot\dfrac{CE}{CA}=\dfrac{ab}{a+b}=8 .
\]
This gives $ab-8a-8b=0$, that is $(a-8)(b-8)=64$. Since $\dfrac{ab}{a+b}<b$, we have $b>8$, so both factors are positive, and $a>b$ means $a-8>b-8$:
\[
(a-8,b-8)\in\{(64,1),(32,2),(16,4)\},\qquad (a,b)\in\{(72,9),(40,10),(24,12)\}.
\]
Each pair is possible with integer legs: an isosceles trapezoid with bases $a,b$ and both legs equal to $a$ exists, since $2a>a-b$ (the drawn one has $a=24$, $b=12$ and legs $10$). So the possible lengths of $AB$ are $72,40,24$, and their sum is $\boxed{136}$.
